\section{Iterate geometry and projective graph classes}\label{sec:geometry}

Fix $n,r\geq1$ and write
\begin{equation}\label{eq:gDQ}
 g=H^n=(g_1,g_2),\qquad D=d^n,\qquad Q=q^r.
\end{equation}
Unless stated otherwise, all varieties, graphs and local rings are over
the algebraically closed field $k$. Use homogeneous coordinates
$[x:y:z]$ on $\Pj^2$, and let $\Linfty=\{z=0\}$.

\begin{lemma}[Degrees and the two infinity points]\label{lem:degrees}
The coordinate polynomials of $g$ satisfy
\[
 \deg g_1=d^{n-1},\qquad \deg g_2=D,
\]
and the degree-$D$ homogeneous part of $g_2$ is $L_ny^D$ for some
$L_n\in\F_q^*$. The degree of $g^{-1}$ is also $D$; its first coordinate
has leading homogeneous part $K_nx^D$, where $K_n\ne0$, and its second
coordinate has degree $d^{n-1}$.

The projective rational extension of $g$ is regular outside
$\Ip=[1:0:0]$. It sends $\Linfty\setminus\{\Ip\}$ to
$\Imn=[0:1:0]$ and fixes $\Imn$. The inverse extension is regular outside
$\Imn$, sends $\Linfty\setminus\{\Imn\}$ to $\Ip$, and fixes $\Ip$.
\end{lemma}

\begin{proof}
Let $c_d\in\F_q^*$ be the leading coefficient of $f$. For $n=1$,
$g_1=y$ and $g_2=f(y)-ax$ give the asserted degrees and $L_1=c_d$.
Suppose the forward assertions hold for $n$. Composition on the left by
$H$ gives
\[
 H^{n+1}=\bigl(g_2,f(g_2)-ag_1\bigr).
\]
The degrees of $f(g_2)$ and $ag_1$ are $d^{n+1}$ and $d^{n-1}$,
respectively. They are unequal, so the leading term cannot cancel.
The leading homogeneous part of the second coordinate is
$c_dL_n^d y^{d^{n+1}}$, and $c_dL_n^d\ne0$. This proves the forward
induction in every characteristic.

For the inverse, the case $n=1$ follows from \eqref{eq:inverse}, with
$K_1=c_d/a$. If $H^{-n}=(u_1,u_2)$ has first degree $d^n$, leading
term $K_nx^{d^n}$, and second degree $d^{n-1}$, then
\[
 H^{-(n+1)}=\bigl((f(u_1)-u_2)/a,u_1\bigr).
\]
Its first degree is $d^{n+1}$, with leading term
$(c_d/a)K_n^d x^{d^{n+1}}$, and its second degree is $d^n$.
This establishes the inverse assertions without a separability condition
on either coordinate polynomial.

Let $G_1,G_2$ be the homogenizations of $g_1,g_2$ to the common degree $D$.
The projective extension is
\begin{equation}\label{eq:homogeneous}
 [x:y:z]\longmapsto[G_1(x,y,z):G_2(x,y,z):z^D].
\end{equation}
Because $\deg g_1<D$, one has $G_1(x,y,0)=0$, whereas
$G_2(x,y,0)=L_ny^D$. There is no indeterminacy on $z\ne0$, since
\eqref{eq:map} is a polynomial morphism. On $z=0$ the three coordinates
in \eqref{eq:homogeneous} vanish simultaneously exactly when $y=0$.
This is the single point $\Ip$. For $y\ne0$ their value is $[0:1:0]$,
which also shows regularity and fixedness at $\Imn$. Homogenizing the
inverse coordinates gives, at $z=0$, the triple $[K_nx^D:0:0]$.
It has base point $\Imn$ and the asserted value $\Ip$ everywhere else
on $\Linfty$.
\end{proof}

Let $X=\Pj^2\times\Pj^2$. The notation $\Gamma_g$ denotes the closure
in $X$ of the graph of $g$ on $\A^2$. The map $\Phi_r$ extends to the
morphism
\[
 [x:y:z]\longmapsto[x^Q:y^Q:z^Q]
\]
on all of $\Pj^2$, so its graph $\Gamma_{\Phi_r}$ is already closed.

\begin{lemma}[Graph charts]\label{lem:graphcharts}
The surface $\Gamma_g$ is the transpose of $\Gamma_{g^{-1}}$.
Over an open subset of the first factor where $g$ is regular it is the
ordinary graph of $g$. Over an open subset of the second factor where
$g^{-1}$ is regular it is parameterized by
$Y\mapsto(g^{-1}(Y),Y)$. In particular, $\Gamma_g$ is smooth at
$(\Imn,\Imn)$ and $(\Ip,\Ip)$.
\end{lemma}

\begin{proof}
On $\A^2\times\A^2$, the graph of an automorphism is the transpose of
the graph of its inverse. Taking closures gives the first assertion.
If $g$ is regular on an open set $U\subset\Pj^2$, its graph over $U$
is closed in $U\times\Pj^2$, since the target is separated. The graph
over $U\cap\A^2$ is dense in that graph, because projection identifies
it with the dense open subset $U\cap\A^2$ of $U$. Thus restriction of
$\Gamma_g$ to the first-factor open set $U$ is exactly that graph.
Applying this reasoning to $g^{-1}$ and transposing proves the second
chart assertion. Lemma~\ref{lem:degrees} supplies the required regular
neighborhoods of $\Imn$ and $\Ip$. Each graph chart is isomorphic to
an open subset of the smooth surface $\Pj^2$.
\end{proof}

Write $h_1$ and $h_2$ for the pullbacks of the hyperplane class from the
two factors of $X$. We use the classical ring and degree convention
\begin{equation}\label{eq:chowring}
 A^*(X)=\Z[h_1,h_2]/(h_1^3,h_2^3),\qquad
 \deg(h_1^2h_2^2)=1.
\end{equation}
The three codimension-two coefficients of a surface in this ring are
determined by intersecting it with two generic lines on one factor or
one generic line on each factor. These are intersection degrees,
so inseparable fibers are counted with their scheme lengths.

\begin{lemma}[Graph classes]\label{lem:classes}
In the convention \eqref{eq:chowring},
\begin{equation}\label{eq:classes}
 [\Gamma_g]=h_1^2+D h_1h_2+h_2^2,
 \qquad
 [\Gamma_{\Phi_r}]=Q^2h_1^2+Qh_1h_2+h_2^2.
\end{equation}
Consequently their Chow intersection has degree
\begin{equation}\label{eq:total}
 \deg\bigl([\Gamma_g]\cdot[\Gamma_{\Phi_r}]\bigr)
 =1+DQ+Q^2.
\end{equation}
\end{lemma}

\begin{proof}
Both projections from $\Gamma_g$ to $\Pj^2$ are birational, as is seen
on the dense affine graph of the automorphism. Thus
$\deg([\Gamma_g]h_1^2)=1$ and
$\deg([\Gamma_g]h_2^2)=1$. These give the coefficients of $h_2^2$
and $h_1^2$ in \eqref{eq:classes}.

For the mixed coefficient choose a line $L$ in the first factor avoiding
$\Ip$, and a generic line $L'$ in the second factor avoiding $\Imn$.
On $L$ the rational extension of $g$ is a morphism. The equation of
$g^{-1}(L')$ is a generic linear combination of $G_1,G_2,z^D$, of degree
$D$. Its restriction to $L\simeq\Pj^1$ is a nonzero section of
$\mathcal O_L(D)$, and hence has zero scheme of length $D$. It is nonzero
at the unique point $L\cap\Linfty$, because that point maps to
$\Imn\notin L'$. The graph chart over $L$ is therefore sufficient to
compute the mixed intersection, which is $D$.

For $\Gamma_{\Phi_r}$ the first projection is an isomorphism, giving
the coefficient $1$ of $h_2^2$. The pullback of a line by $\Phi_r$ has
degree $Q$, giving the coefficient $Q$ of $h_1h_2$. The second projection
is the finite morphism $\Phi_r$, of degree $Q^2$. For example, its
function-field extension on the affine chart is
\[
 k(x^Q,y^Q)\subset k(x,y),
\]
and the $Q^2$ monomials $x^iy^j$, $0\leq i,j<Q$, form a basis of the
larger field over the smaller one. This yields the coefficient $Q^2$
of $h_1^2$. Here $Q^2$ is the degree of a finite morphism, not the
cardinality of a geometric fiber: the Frobenius morphism is radicial
and has one point in each geometric fiber.

Multiplying the two classes, the only surviving degree-four monomial
is $h_1^2h_2^2$. Its coefficient is $1+DQ+Q^2$, proving
\eqref{eq:total}. We have computed a Chow product; the next section
checks properness before using it as a sum of local intersection numbers.
\end{proof}
